% Ampliación para insertar después del bloque térmico de hibridacion.tex.
% Procedencia, dependencias y alcance de los controles: INFORME.md.
\section{Lectura marcada y conjugación termométrica de la acción}
\label{md:r27:termica}

La acción orientada, el calendario, los lectores de capacidad y el archivo de
memoria proceden de la misma dinámica enriquecida APP--TRIT--TPK. Las dos hojas
conservan residuo y cociente; el régimen trítico conserva orientación; el
transporte $U_t$ y su holonomía $\Gamma_9$ preceden a la lectura reducida
$K_{\rm ph}$. El retorno de fase conserva una historia que sigue creciendo.
La presente construcción compone esos resultados en un registro de referencia
y ensayo, dentro de la estructura discreta conjunta del continuo ya construida.
Su regla termométrica se especifica mediante dos operaciones sobre ese registro.

\subsection{Capacidad, referencia espectral y origen retenido}
\label{md:r27:capacidad}
La capacidad de un prefijo de índice $t\ge0$ y su resto son
\begin{equation}
 C_t=\lfloor\log_{10}(9^{t+1})\rfloor,\qquad
 9^{t+1}=10^{C_t}R_t,\qquad 1\le R_t<10.
 \label{md:r27:capacidad-resto}
\end{equation}
El conteo se evalúa por potencias enteras. Las desigualdades
$10^{20}\le9^{21}<10^{21}$ y $10^{20}\le9^{22}<10^{21}$,
junto con $10^{j-1}\le9^j<10^j$ para $1\le j\le21$, prueban que
la primera vacancia positiva tiene índice $21$ y capacidad $20$.
En la lectura trítica de capacidad, los residuos $0,1,2$ corresponden a
$+1,-1,0$, respectivamente. Las capacidades en $t=21,22,23,24$ son
$20,21,22,23$; el primer retorno posterior del régimen es, por tanto,
$t=24$. El prefijo contiene $25$ posiciones y expresa capacidad $23$.
Así se conservan las funciones diferenciadas de $(20,24,25,23)$.
El regreso del régimen conserva una fase nonádica distinta; la fase retorna
en $t=30$, con capacidad $29$.

El portador del lector es $\mathcal D=\mathcal D_0\oplus\mathcal D_1
\oplus\mathcal D_2$, con rangos $(1,380,649)$ y graduación $N=0,1,2$.
La construcción conserva dos graduaciones: $N$ distingue el prefijo y $L$
la ocupación del espacio exterior o simétrico. Para $m=20$, $n=24$,
sean $\mathcal F_2^J$ el fijo exterior de grado dos de dos hojas intercambiadas
y $\mathcal B_{\le2}^J$ el fijo simétrico de grado a lo sumo dos. El portador es
$\mathbb C\Omega_{\rm ext}\oplus\mathcal F_2^J\oplus\mathcal B_{\le2}^J$.
La involución $J$ intercambia las dos hojas mediante el intercambio graduado;
en el sector exterior cada uno de sus autoespacios tiene dimensión $m$.
Los caracteres de los subespacios fijos son
\begin{equation}
 \begin{aligned}
 \chi_F(z)&=\frac{(1+z)^{2m}+(1-z^2)^m}{2},\\
 \chi_B(z)&=\frac{P_n(z)^2+P_n(z^2)}{2},\\
 P_n(z)&=\prod_{j\ge1}(1-z^j)^{-n}.
 \end{aligned}
 \label{md:r27:caracteres-fock}
\end{equation}
En efecto, el proyector $(I+J)/2$ promedia el carácter total y la traza
del intercambio. En el espacio exterior, los autovalores $+1,-1$ de $J$
producen $(1+z)^m(1-z)^m$; en el producto simétrico, sólo los pares idénticos
contribuyen a la traza del intercambio, produciendo $P_n(z^2)$.
Para los coeficientes hasta grado dos basta el polinomio
$P_n(z)=1+nz+n(n+3)z^2/2+O(z^3)$; los productos formales son finitos en
cada grado. Extraer los coeficientes demuestra
\begin{equation}
 \begin{split}
 \dim\mathcal F_2^J&=m(m-1)=380,\\
 (\dim\mathcal B_0^J,\dim\mathcal B_1^J,\dim\mathcal B_2^J)
 &=(1,n,n(n+2))=(1,24,624).
 \end{split}
 \label{md:r27:grados-fock}
\end{equation}
Por tanto, la descomposición simultánea $(N,L;\text{multiplicidad})$ es
\[
 \begin{gathered}
 (0,0;1),(1,2;380),(2,0;1),\\
 (2,1;24),(2,2;624).
 \end{gathered}
\]
El rango $649=1+24+624=25^2+24$ conserva tres ocupaciones; el vacío externo
y el vacío del sumando simétrico permanecen ortogonales. El carácter conjunto
es $1+380xy^2+x^2(1+24y+624y^2)$.
La asignación de $N=0,1,2$ a los cortes $24,27,30$, de capacidades
$23,26,29$, especifica el lector utilizado en esta realización; mantiene $L$
como observable concomitante. En particular,
$\operatorname{Tr}q^L=2+24q+1004q^2$ y
$\operatorname{Tr}q^N=1+380q+649q^2$ conservan operaciones distintas.

El origen de la referencia espectral reside en
\eqref{md:antmem:cuantica:eq:5.1}--\eqref{md:antmem:cuantica:eq:5.10}.
El acoplamiento nonádico $8:1$ y el lazo ordenado
$H_1=\left(\begin{smallmatrix}2&1\\1&1\end{smallmatrix}\right)$,
con la preparación de dos modos puros idénticos y la coordenatización de covarianza
transportada conjuntamente, dan
$\nu^2=1+(8/81)[\operatorname{tr}(H_1H_1^{\mathsf T})-2]=121/81$.
La suma de cuadrados de las entradas de $H_1$ es siete, de modo que
$\nu=11/9$. El teorema~\ref{md:antmem:gauss:thm:espectro-gaussiano}
da entonces $r=(\nu-1)/(\nu+1)=1/10$ y autovalores $(1-r)r^j$.
Se utiliza esa preparación demostrada, incluida su reducción de memoria.
La división $j=3n+b$, $0\le b<3$, conserva cociente y residuo y satisface
\begin{equation}
 (1-r)r^{3n+b}=(1-r^3)(r^3)^n\frac{r^b}{1+r+r^2}.
 \label{md:r27:factorizacion}
\end{equation}
Por consiguiente, el cociente tiene razón $q=r^3=1/1000$ y el residuo conserva
la distribución $\operatorname{diag}(100,10,1)/111$. El aumento del cociente
tras completar los tres residuos conserva el acarreo. Si la energía elemental
es $\epsilon$, la energía del índice es $\epsilon(3n+b+1/2)$: el salto entre
bloques es $3\epsilon$ a la misma temperatura de la referencia.

Definimos el operador positivo de referencia y su intensidad:
\begin{equation}
 \begin{gathered}
 \eta=\bigoplus_{n=0}^{2}r^{23+3n}I_{d_n},\qquad
 (d_0,d_1,d_2)=(1,380,649),\\
 b_*:=\operatorname{Tr}\eta=\frac{1380649}{10^{29}}.
 \end{gathered}
 \label{md:r27:eta}
\end{equation}
En particular, $b_*=1{,}380649\times10^{-23}$ y $0<b_*<1$.
El origen $23$ está retenido en $\eta$; normalizar únicamente su estado
elimina ese factor de las probabilidades relativas.

La sección de acción de \eqref{md:hib:reloj} proporciona
$\epsilon_a=h_a/(108t_0)$ y $\beta_{\rm clk}=\ln3/\epsilon_a$.
Para la elevación energética aditiva de este lector de capacidad,
\begin{equation}
 \begin{split}
 H_{\rm tot}(t)&=2(t+1)\epsilon_a,\\
 H_{\rm mem}(t)&=\frac{\epsilon_a}{\ln3}\ln R_t,\\
 H_{\rm cap}(t)&=H_{\rm tot}(t)-H_{\rm mem}(t)
 =\frac{\epsilon_a\ln10}{\ln3}C_t.
 \end{split}
 \label{md:r27:energia-capacidad}
\end{equation}
La prueba consiste en tomar logaritmos en \eqref{md:r27:capacidad-resto}.
Por sustitución, $e^{-\beta_{\rm clk}H_{\rm cap}(t)}=10^{-C_t}$.
Sobre $\mathcal D$, esto da
$H_{\rm cap}=(\epsilon_a\ln10/\ln3)(23I+3N)$ y
$\operatorname{Tr}e^{-\beta_{\rm clk}H_{\rm cap}}=b_*$.
El resto permanece registrado como contribución energética; la compensación
es una lectura distinta de eliminar un factor por traza parcial.

\subsection{Identidad tensorial y especificación de los lectores}
\label{md:r27:lectores}
\begin{lemma}[Información incremental del registro producto]
Para toda matriz positiva $\eta$ y toda matriz de densidad $\rho$, con
$\mathscr H(X)=-\operatorname{Tr}(X\ln X)$ y $0\ln0=0$, se cumple
\begin{equation}
 \begin{gathered}
 \mathscr H(\eta\otimes\rho)-\mathscr H(\eta)
 =\operatorname{Tr}(\eta)s(\rho),\\
 s(\rho)=-\operatorname{Tr}(\rho\ln\rho).
 \end{gathered}
 \label{md:r27:tensor}
\end{equation}
\end{lemma}
\begin{proof}
Sobre los soportes, el logaritmo del producto tensorial es
$\ln\eta\otimes I+I\otimes\ln\rho$. La factorización de la traza y
$\operatorname{Tr}\rho=1$ dan
$\mathscr H(\eta\otimes\rho)=\mathscr H(\eta)+\operatorname{Tr}(\eta)s(\rho)$.
La continuidad de $x\ln x$ en cero extiende la igualdad a núcleos no triviales.
\end{proof}

Para \eqref{md:r27:eta}, formamos un registro normalizado y su marca:
\begin{equation}
 \widehat\eta=\eta\oplus(1-b_*)|f\rangle\langle f|,
 \qquad\Omega_\rho=\widehat\eta\otimes\rho,\qquad
 P=I_{\mathcal D}\oplus0,
 \label{md:r27:registro}
\end{equation}
donde $P$ actúa por la identidad sobre el factor de ensayo. La bandera $f$
completa la probabilidad de esta realización. Cualquier archivo microscópico
adicional se transporta como tal; la bandera conserva exclusivamente esta
distinción de selección. Para un Hamiltoniano fijo $H$, se especifican
\begin{equation}
 \begin{split}
 \mathcal S_P(\rho)&=-\operatorname{Tr}
 [P\Omega_\rho(I\otimes\ln\rho)]=b_*s(\rho),\\
 \mathcal E_P(\rho)&=
 \frac{\operatorname{Tr}[P\Omega_\rho(I\otimes H)]}
 {\operatorname{Tr}(P\Omega_\rho)}=\operatorname{Tr}(\rho H).
 \end{split}
 \label{md:r27:par-lector}
\end{equation}
La primera operación lee información antes del condicionamiento; la segunda
lee energía condicionada a la marca. Ambas operan sobre el mismo registro,
y la probabilidad de la marca sigue siendo $b_*$. La referencia permanece
fija al variar el ensayo.

\begin{theorem}[Conjugación termométrica del lector marcado]
Sea $H$ un Hamiltoniano autoadjunto fijo sobre un espacio finito y sea
$\rho_\beta=e^{-\beta H}/Z(\beta)$, con $\beta>0$ y
$\operatorname{Var}_{\rho_\beta}(H)>0$. Para el par
\eqref{md:r27:par-lector}, la temperatura conjugada satisface
\begin{equation}
 \Theta_P:=\left(\frac{d\mathcal S_P}{d\mathcal E_P}\right)^{-1},
 \qquad b_*\Theta_P=\beta^{-1}.
 \label{md:r27:transductor}
\end{equation}
\end{theorem}
\begin{proof}
Escribamos $E(\beta)=\operatorname{Tr}(\rho_\beta H)$.
Derivar la suma espectral finita da
$Z'/Z=-E$ y $E'=-\operatorname{Var}_{\rho_\beta}(H)$.
Como $s(\rho_\beta)=\beta E+\ln Z$, se tiene $s'=\beta E'$.
La referencia fija permite derivar $\mathcal S_P=b_*s$ sin término adicional,
y la varianza positiva permite dividir por $E'$. Por tanto,
$d\mathcal S_P/d\mathcal E_P=b_*\beta$, que prueba el enunciado.
\end{proof}

La normalización es parte de la definición operativa. Si información y energía
se leen ambas condicionadas, sus valores son $(s,E)$ y su razón diferencial es
$\beta$. Si ambas se leen antes de condicionar, sus valores son $(b_*s,b_*E)$
y la misma razón es $\beta$. El par \eqref{md:r27:par-lector} produce
$b_*\beta$ porque conserva la intensidad en la primera operación y utiliza
energía por ensayo seleccionado en la segunda. Mantener ese par fija su
coeficiente: reemplazarlo por $cb_*$, con $c\ne1$, cambia la ecuación.
Si la referencia variase, $d(b_*s)=b_*\,ds+s\,db_*$; ese protocolo tendría
otra respuesta y exige conservar el segundo término.

\subsection{Selección variacional y balance de información}
\begin{theorem}[Principio variacional en la realización marcada]
Para $\Theta>0$, el funcional
\begin{equation}
 \Phi_\Theta(\rho)=\mathcal E_P(\rho)-\Theta\mathcal S_P(\rho)
 =\operatorname{Tr}(\rho H)-b_*\Theta s(\rho)
 \label{md:r27:funcional}
\end{equation}
tiene un único mínimo sobre las matrices de densidad, dado por
\begin{equation}
 \rho_\Theta^*=
 \frac{e^{-H/(b_*\Theta)}}{\operatorname{Tr}e^{-H/(b_*\Theta)}}.
 \label{md:r27:gibbs}
\end{equation}
Además,
\begin{equation}
 \Phi_\Theta(\rho)-\Phi_\Theta(\rho_\Theta^*)
 =b_*\Theta D(\rho\Vert\rho_\Theta^*)\ge0.
 \label{md:r27:variacional}
\end{equation}
\end{theorem}
\begin{proof}
El estado \eqref{md:r27:gibbs} es fiel en dimensión finita.
Sustituir $\ln\rho_\Theta^*=-H/(b_*\Theta)-\ln Z_\Theta$ en
$D(\rho\Vert\rho_\Theta^*)=\operatorname{Tr}\rho(\ln\rho-\ln\rho_\Theta^*)$
da la igualdad. Para verificar la positividad, diagonalícense $\rho$ y el
estado fiel $\sigma=\rho_\Theta^*$, con autovalores $r_i,s_j$ y matriz
bistocástica de solapamientos $p_{ij}$. La concavidad de $\ln$ da
$D(\rho\Vert\sigma)\ge\sum_i r_i\ln(r_i/t_i)$,
donde $t_i=\sum_jp_{ij}s_j>0$ y $\sum_i t_i=1$.
La desigualdad $-\ln u\ge1-u$ demuestra que esta última suma es no negativa.
La igualdad exige $r_i=t_i$ para todos los índices y que cada vector propio
de $\rho$ esté contenido en un autoespacio de $\sigma$ de valor $t_i$;
por tanto, exige $\rho=\sigma$. Esto prueba el mínimo y su unicidad.
La convención sobre el núcleo incluye estados singulares; la fidelidad de
$\sigma$ impide que uno de ellos alcance la igualdad.
\end{proof}
En esta realización, el mismo coeficiente interviene en la conjugación
diferencial y en el equilibrio variacional. La generalización a sistemas
infinitos conserva por separado las condiciones de traza, diferenciabilidad
y finitud de energía y entropía; los teoremas anteriores tienen dominio finito.

\subsection{Composición con acción, área y reciprocidad gravitatoria}
\label{md:r27:composicion}
Se utilizan las secciones orientadas ya construidas, en una misma coordenatización:
\begin{equation}
 \begin{gathered}
 \epsilon_a=\frac{h_a}{108t_0}=\frac{\hbar_a c_{\rm int}}{L_\alpha},
 \qquad q_A=\gamma a_0L_\alpha^2,\\
 \mathcal T_{I/A}=\frac{\epsilon_a}{q_A},\qquad
 G_a=\frac{c_{\rm int}^3L_\alpha^2}{\hbar_a}.
 \end{gathered}
 \label{md:r27:secciones}
\end{equation}
La igualdad longitudinal conserva el reloj de
\eqref{md:rigidez:reloj}; $q_A>0$ fija la orientación areal utilizada.
El funcional $\gamma$ es el de \eqref{md:compos:funcional-barbero} y mantiene
las incidencias hexada--octada y los canales $90/120$. Su aparición en el área
conserva esa dependencia completa. La sección~\ref{md:rigidez:area}
construye $a_0$ como carácter trítico angular normalizado y demuestra
$\mathcal A=q_A(N_{90}+N_{120})$ sobre treinta y seis enlaces etiquetados.
Así, $q_A$ es el incremento de área por unidad de ocupación en esa sección,
y no un factor inferido de una entropía prefijada.

\begin{corollary}[Compatibilidad de los lectores térmico, areal y gravitatorio]
Con la referencia trítica $\beta_{\rm clk}=\ln3/\epsilon_a$ y el mismo par
\eqref{md:r27:par-lector}, se cumple
\begin{equation}
 \begin{gathered}
 b_*\Theta_{{\rm clk},P}
 =\frac{\epsilon_a}{\ln3}
 =\frac{\mathcal T_{I/A}\gamma a_0L_\alpha^2}{\ln3},\\
 G_a b_*\Theta_{{\rm clk},P}\ln3=c_{\rm int}^4L_\alpha.
 \end{gathered}
 \label{md:r27:area-gravedad}
\end{equation}
Para la sección de horizonte $\tau_R=\hbar_a c_{\rm int}/(2\pi R)$,
\begin{equation}
 b_*T_{R,P}=\tau_R,\qquad
 \frac{T_{R,P}}{\Theta_{{\rm clk},P}}
 =\frac{\ln3}{2\pi}\frac{L_\alpha}{R}.
 \label{md:r27:horizonte}
\end{equation}
\end{corollary}
\begin{proof}
Aplicar \eqref{md:r27:transductor} al reloj da la primera igualdad.
Sustituir $q_A$ en $\mathcal T_{I/A}q_A=\epsilon_a$ da la segunda.
Multiplicar por $G_a$ y usar \eqref{md:r27:secciones} da
$G_a\epsilon_a=c_{\rm int}^4L_\alpha$. Para el horizonte se aplica el mismo
teorema con $\beta_R=1/\tau_R$ y se divide por la igualdad del reloj.
\end{proof}
Las conversiones entre secciones $i,j$ conservan
$\Theta_{i,P}/\Theta_{j,P}=\tau_i/\tau_j$. Todo ciclo de conversiones
se cierra por cancelación telescópica. La temperatura del horizonte y la del
reloj coinciden sólo cuando coinciden sus secciones energéticas.

En una coordenatización interna, $b_*$ es el coeficiente del transductor definido por
\eqref{md:r27:par-lector}. La unidad térmica se especifica mediante la
conjugación entre esas lecturas y la sección energética elegida. La comparación
con la aplicación denominada Boltzmann en otro protocolo transporta también
sus lectores, sus rectas y sus bases; la cifra sola conserva el alcance escalar.
En particular, esta ampliación introduce una realización termométrica
explícita y sus pruebas. Su regla se incorpora como tal, sin atribuirla
retrospectivamente a todo operador térmico anterior ni identificar por nombre
su unidad interna con una unidad metrológica externa.
